\documentclass[12pt]{book}
\usepackage{precalc1}
\setcounter{chapter}{0}
\begin{document}

\chapter{Functions}

A relation is a collection of points. A function is a relation in which each input
determines a single output, which is precisely the condition that lets you predict.
Data arrives as a scatter of points, and the work of the subject is to find a
predictable rule running through it: a function fit to the cloud, reshaped from a
known curve through $A\,f\!\big(b(x-h)\big)+k$.

\section{Relations as Point Clouds}

Record, for every student in a class, their height and their shoe size. Each
student becomes a pair of numbers: an input (height) and an output (shoe size).
Plotted, these pairs form a cloud of dots. The cloud is the raw object of study. It
carries no formula and no curve; it is the data, exactly as measured.

\begin{definition}{Relation}
A \emph{relation} is any set of ordered pairs $(x,y)$. The set of all first
coordinates is the \emph{domain}; the set of all second coordinates is the
\emph{range}.
\end{definition}

A relation can be specified by a list, a table, an equation, or a graph. These are
four ways of naming the same set of pairs. The list
\[
\{(1,2),\ (3,5),\ (3,-1),\ (4,0)\}
\]
is a relation with domain $\{1,3,4\}$ and range $\{2,5,-1,0\}$. So is the set of all
$(x,y)$ satisfying $x^2+y^2=1$; that relation is a circle, an infinite set of points
packed densely enough to read as a curve.

\begin{visual}
A relation is a cloud of dots suspended above the $x$-axis. The domain is the shadow
the cloud casts straight down onto the $x$-axis; the range is the shadow it casts
sideways onto the $y$-axis. Domain and range are the two shadows of the cloud.
\end{visual}

\begin{center}
\begin{tikzpicture}[scale=0.85]
  \draw[->] (-0.5,0) -- (6,0) node[right] {$x$};
  \draw[->] (0,-0.5) -- (0,4.5) node[above] {$y$};
  % point cloud
  \foreach \p in {(1,1.2),(1.3,1.8),(1.1,2.4),(2,2.1),(2.4,2.9),(2.1,1.6),
                  (3,3.1),(3.3,2.6),(3.1,3.6),(4,3.4),(4.4,3.9),(4.1,2.9),
                  (1.8,1.0),(2.7,2.3),(3.7,3.2),(4.6,3.5)}{
    \fill[cloudgray] \p circle (2pt);
  }
  % domain shadow
  \draw[pcblue,line width=2pt] (1,-0.15) -- (4.6,-0.15);
  \node[pcblue] at (5.6,-0.15) {\small domain};
  % range shadow
  \draw[pcgreen,line width=2pt] (-0.15,1.0) -- (-0.15,3.9);
  \node[pcgreen,rotate=90] at (-0.5,2.5) {\small range};
\end{tikzpicture}
\end{center}

\begin{example}[reading domain and range from a list]
For $\{(-2,4),(0,0),(1,1),(2,4)\}$ the domain is $\{-2,0,1,2\}$ and the range is
$\{0,1,4\}$. The value $4$ appears twice in the range but is listed once; a range
is a set, so repeats collapse.
\end{example}

\section{When a Relation Is a Function}

Most clouds are unstructured. A few satisfy a strong condition: each input is paired
with exactly one output. That condition is what makes prediction possible. Given the
input, there is one permitted output, so the output can be stated. If two outputs
share an input, prediction fails, because the input no longer selects a single
value.

\begin{definition}{Function}
A \emph{function} is a relation in which each input is paired with exactly one
output. Equivalently, no two distinct pairs share the same first coordinate.
\end{definition}

The relation $\{(1,2),(3,5),(3,-1),(4,0)\}$ is not a function: the input $3$ is
paired with both $5$ and $-1$. Deleting either pair makes it a function. The circle
$x^2+y^2=1$ is not a function either; the input $x=0$ yields both $y=1$ and $y=-1$.

\begin{keyidea}
A function assigns to each input exactly one output. The domain is the set of inputs
for which the assignment is defined; for each such input the output is determined.
\end{keyidea}

When a function is given by a rule, $f(x)$ denotes the output at input $x$. The
symbol $f(2)$ means the output of $f$ at input $2$, a single number. It does not mean
$f$ multiplied by $2$.

\begin{example}[evaluating a function]
Let $f(x)=x^2-3x+1$. Then
\[
f(2)=4-6+1=-1,\qquad f(0)=1,\qquad f(-1)=1+3+1=5.
\]
And $f(a+1)=(a+1)^2-3(a+1)+1=a^2-a-1$ after expanding. Evaluating at an expression
means substituting the entire expression for $x$.
\end{example}

\begin{pitfall}
$f(a+b)$ is almost never $f(a)+f(b)$. For $f(x)=x^2$, $f(1+2)=f(3)=9$, while
$f(1)+f(2)=1+4=5$. The function acts on the entire input at once, not on its parts
separately.
\end{pitfall}

\subsection{The Vertical Line Test}

The function property has a graphical reading. ``Each input gives one output'' means
that above any point on the $x$-axis the graph has at most one point. A vertical line
is the set of points sharing one $x$-value, which gives the test directly.

\begin{keyidea}
\emph{Vertical line test.} A graph represents a function if and only if no vertical
line crosses it more than once.
\end{keyidea}

\begin{center}
\begin{tikzpicture}[scale=0.8]
  % parabola opening up: a function
  \begin{scope}
  \draw[->] (-2.2,0) -- (2.2,0) node[right] {$x$};
  \draw[->] (0,-0.4) -- (0,3.2) node[above] {$y$};
  \draw[pcblue,thick,domain=-1.6:1.6,smooth] plot (\x,{\x*\x});
  \draw[pcred,dashed] (0.8,-0.2) -- (0.8,3.1);
  \fill[pcred] (0.8,0.64) circle (1.6pt);
  \node at (0,-1) {\small passes: a function};
  \end{scope}
  % sideways parabola: not a function
  \begin{scope}[xshift=6cm]
  \draw[->] (-0.4,0) -- (3.4,0) node[right] {$x$};
  \draw[->] (0,-1.8) -- (0,1.8) node[above] {$y$};
  \draw[pcblue,thick,domain=-1.6:1.6,smooth] plot ({\x*\x},\x);
  \draw[pcred,dashed] (1.4,-1.7) -- (1.4,1.7);
  \fill[pcred] (1.4,1.183) circle (1.6pt);
  \fill[pcred] (1.4,-1.183) circle (1.6pt);
  \node at (1.5,-2.6) {\small fails: not a function};
  \end{scope}
\end{tikzpicture}
\end{center}

\begin{example}[applying the test]
The graph of $y=|x|$ is a V opening upward; every vertical line hits it once, so it
is a function. The graph of $x=y^2$ is a sideways parabola; the vertical line
$x=1$ hits it at $(1,1)$ and $(1,-1)$, so it is not a function.
\end{example}

\section{Fitting a Function to a Cloud}

The two ideas now meet. Real data forms a point cloud with noise; it is almost never
an exact function. A function can nonetheless be laid through the cloud as a
\emph{model}: a predictable rule that follows the trend closely enough to forecast,
summarize, and reason about the data.

\begin{keyidea}
A point cloud is the data. A function fit through it is a \emph{model}. The model is
not the data; it is the predictable rule chosen to represent the data so that one can
compute and forecast with it.
\end{keyidea}

The first decision is the \emph{shape}. A cloud rising in a straight band calls for a
linear model. A cloud curving upward like a bowl suggests a parabola. A cloud that
climbs quickly and then levels off suggests a square-root or logarithmic shape. The
catalog of basic shapes is the library of \emph{parent functions}, and each can be
reshaped to fit by the four operations that follow.

\demo{A hidden transformed parent function generates a noisy point cloud. Identify
the parent-function family, then tune $y=A\,f(b(x-h))+k$ until the curve threads the
cloud. The fit score reports the root-mean-squared error between curve and data; the
hidden answer is deliberately approximate because the cloud carries
noise.}{https://bobovski66.github.io/Precalculus/demos/function\_fit.html}

\begin{example}[choosing a shape]
A cloud trends upward and bends, steep at first and then flatter. A straight line
undershoots the early rise and overshoots the later leveling. The square-root shape
$f(x)=\sqrt{x}$ has exactly this profile, fast and then flattening, so $\sqrt{x}$ is
the parent to start from; the fitting work is sliding and stretching it onto the
cloud.
\end{example}

\section{Transformations: Bending One Curve into Another}

Every function in the library can be reshaped by four operations, all captured in a
single expression. Given a parent function $f$, form
\[
g(x)=A\,f\!\big(b(x-h)\big)+k.
\]
Each constant controls one geometric action. Fitting reduces to reading a cloud and
setting these four numbers.

\begin{keyidea}
\textbf{The four constants of} $\;g(x)=A\,f\!\big(b(x-h)\big)+k$.
\begin{itemize}\setlength{\itemsep}{2pt}
  \item $k$: vertical shift. Adds to the output, moving the graph up ($k>0$) or
        down ($k<0$).
  \item $h$: horizontal shift. Moves the graph right ($h>0$) or left ($h<0$).
  \item $A$: vertical stretch by factor $|A|$; if $A<0$, also a flip across the
        $x$-axis.
  \item $b$: horizontal stretch by factor $1/|b|$; if $b<0$, also a flip across the
        $y$-axis.
\end{itemize}
\end{keyidea}

\subsection{Outside versus inside}

The reliable bookkeeping is the outside/inside rule. Constants applied \emph{outside}
$f$ ($A$ and $k$) act on the output, so they move the graph \emph{vertically} and
behave as expected: $+k$ raises it, $A$ scales its height. Constants applied
\emph{inside} $f$ ($b$ and $h$) act on the input, so they move the graph
\emph{horizontally} and behave in reverse: $x-h$ shifts \emph{right} by $h$, and
multiplying the input by $b$ \emph{compresses} by $b$.

\begin{visual}
Outside changes act on the picture after it is drawn: raise it, scale its height.
Inside changes act on the input before $f$ receives it, so they warp the horizontal
axis, and warping the input runs opposite to intuition. The parent produces its
value at input $0$; the transformed graph reaches that value where $b(x-h)=0$, that
is, at $x=h$. This is why $x-h$ moves the graph right rather than left.
\end{visual}

\begin{pitfall}
The horizontal shift is $x-h$, so $g(x)=f(x-3)$ moves the graph \emph{right} by 3,
not left. The minus sign is built into the form. Read $f(x+3)$ as $f(x-(-3))$, so
$h=-3$ and the shift is left by 3.
\end{pitfall}

\begin{example}[a single vertical shift]
With parent $f(x)=x^2$, the function $g(x)=x^2+3$ is the parabola lifted up 3
units. The vertex moves from $(0,0)$ to $(0,3)$; every output gains 3.
\end{example}

\begin{example}[a single horizontal shift]
With $f(x)=x^2$, the function $g(x)=(x-2)^2$ slides the parabola right 2 units. The
vertex moves from $(0,0)$ to $(2,0)$. Check: $g(2)=0$, the old value at the old
vertex input $0$.
\end{example}

\begin{example}[stretch and reflect together]
With $f(x)=x^2$, the function $g(x)=-2x^2$ has $A=-2$: stretch vertically by 2
(twice as tall) and flip across the $x$-axis (opens downward). The point $(1,1)$ on
the parent becomes $(1,-2)$.
\end{example}

\begin{example}[all four at once]
Let $f(x)=\sqrt{x}$ and
\[
g(x)=3\sqrt{2(x-1)}+4.
\]
Read off $A=3,\ b=2,\ h=1,\ k=4$. Starting from $\sqrt{x}$: compress horizontally by
$\tfrac12$, stretch vertically by 3, shift right 1, shift up 4. The parent begins at
$(0,0)$; the transformed curve begins at $(h,k)=(1,4)$, since the inside is zero at
$x=1$ and the $+k$ raises the value to $4$.
\end{example}

\begin{center}
\begin{tikzpicture}[scale=0.8]
  \draw[->] (-0.5,0) -- (6,0) node[right] {$x$};
  \draw[->] (0,-0.5) -- (0,7) node[above] {$y$};
  \draw[pcblue,thick,domain=0:5.5,smooth] plot (\x,{sqrt(\x)});
  \node[pcblue] at (5.2,1.9) {\small $f(x)=\sqrt{x}$};
  \draw[pcred,thick,domain=1:5.5,smooth] plot (\x,{3*sqrt(2*(\x-1))+4});
  \node[pcred] at (3.6,6.6) {\small $g(x)=3\sqrt{2(x-1)}+4$};
  \fill[pcred] (1,4) circle (2pt) node[below right] {\small $(1,4)$};
  \fill[pcblue] (0,0) circle (2pt) node[below right] {\small $(0,0)$};
\end{tikzpicture}
\end{center}

\subsection{Order of operations on the graph}

When several transformations combine, the correct order mirrors the evaluation of
$g(x)$ at a number: work from the input outward. Apply the inside first (horizontal:
the $b$ compression, then the $h$ shift), then the outside (vertical: the $A$
stretch or reflection, then the $k$ shift). Following one landmark point of the
parent through each step keeps the order honest.

\begin{example}[tracking a point]
Take $f(x)=\sqrt{x}$ and the landmark $(4,2)$ on it. Apply
$g(x)=3\sqrt{2(x-1)}+4$ to that point's coordinates:
\[
x:\quad 4 \xrightarrow{\ \div b\ } 2 \xrightarrow{\ +h\ } 3,
\qquad
y:\quad 2 \xrightarrow{\ \times A\ } 6 \xrightarrow{\ +k\ } 10.
\]
So $(4,2)$ lands at $(3,10)$. Confirm directly: $g(3)=3\sqrt{2\cdot 2}+4
=3\cdot 2+4=10$. The two routes agree.
\end{example}

\begin{pitfall}
The input coordinate transforms by the \emph{inverse} operations in \emph{reverse}
order: to locate where an $x$ lands, undo the inside. Dividing by $b$ before adding
$h$ matches the algebra of $b(x-h)$; reversing those two steps misplaces the point.
\end{pitfall}

\section{The Library of Parent Functions}

Fitting is limited by the shapes one knows. The following parents recur throughout
the course; their graphs should be memorized well enough to recognize in a cloud.

\begin{center}
\renewcommand{\arraystretch}{1.4}
\begin{tabular}{@{}lll@{}}
\toprule
\textbf{Name} & \textbf{Rule} & \textbf{Shape in a sentence}\\
\midrule
Constant     & $f(x)=c$        & flat horizontal line\\
Linear       & $f(x)=x$        & straight line through the origin, slope 1\\
Absolute value & $f(x)=|x|$    & V opening up, corner at origin\\
Quadratic    & $f(x)=x^2$      & bowl opening up, vertex at origin\\
Cubic        & $f(x)=x^3$      & S-curve, flat at origin\\
Square root  & $f(x)=\sqrt{x}$ & half-parabola, fast then flattening\\
Reciprocal   & $f(x)=1/x$      & two branches, axes as asymptotes\\
Exponential  & $f(x)=2^x$      & flat then explosive growth\\
\bottomrule
\end{tabular}
\end{center}

\begin{center}
\begin{tikzpicture}[scale=0.62]
  \begin{scope}
    \draw[->] (-2.2,0)--(2.2,0); \draw[->] (0,-0.4)--(0,4);
    \draw[pcblue,thick,domain=-1.9:1.9,smooth] plot (\x,{\x*\x});
    \node at (0,-1) {\small $x^2$};
  \end{scope}
  \begin{scope}[xshift=5cm]
    \draw[->] (-2.2,0)--(2.2,0); \draw[->] (0,-2)--(0,2);
    \draw[pcblue,thick,domain=-1.25:1.25,smooth] plot (\x,{\x*\x*\x});
    \node at (0,-2.6) {\small $x^3$};
  \end{scope}
  \begin{scope}[xshift=10cm]
    \draw[->] (-0.4,0)--(4.2,0); \draw[->] (0,-0.4)--(0,2.4);
    \draw[pcblue,thick,domain=0:4,smooth] plot (\x,{sqrt(\x)});
    \node at (2,-1) {\small $\sqrt{x}$};
  \end{scope}
  \begin{scope}[xshift=15cm]
    \draw[->] (-2.2,0)--(2.2,0); \draw[->] (0,-0.4)--(0,3.6);
    \draw[pcblue,thick,domain=-1.8:1.8,smooth] plot (\x,{2^\x});
    \node at (0,-1) {\small $2^x$};
  \end{scope}
\end{tikzpicture}
\end{center}

\begin{example}[matching a cloud to a parent]
A cloud lies entirely in the first quadrant, rising steeply near zero and then
bending to a gentle climb. Among the parents, only $\sqrt{x}$ has this steep-then-
gentle profile beginning at the origin. Start from $\sqrt{x}$ and fit $A,b,h,k$.
\end{example}

\section{Exercises}

\begin{exercises}
\item For the relation $\{(-1,3),(0,0),(2,3),(2,-3),(5,1)\}$, state the domain and
range, and decide whether it is a function. Justify.

\item Decide whether each graph could pass the vertical line test, and say what the
test is checking: (a) a circle; (b) a downward parabola; (c) the letter ``S'' lying
on its back; (d) a single steep upward line.

\item Let $f(x)=2x^2-x+3$. Compute $f(0)$, $f(-2)$, and $f(t+1)$, simplifying the
last completely.

\item Give a relation that is \emph{not} a function, then change exactly one ordered
pair to make it a function. Explain why your change works.

\item Starting from $f(x)=x^2$, write the rule for the function obtained by shifting
left 4, stretching vertically by 3, and shifting down 1. Then state the vertex.

\item Describe in words every transformation taking $f(x)=|x|$ to
$g(x)=-\tfrac12\,|x+2|+5$, in the correct order, and give the coordinates of the
corner of $g$.

\item The point $(9,3)$ lies on $f(x)=\sqrt{x}$. Find where it lands under
$g(x)=2\sqrt{x-1}+4$ by tracking the coordinate through each step, then confirm by
direct substitution.

\item Explain, using the outside/inside rule, why $f(x-3)$ shifts the graph right
while $f(x)+3$ shifts it up. Why do the two ``$+3$'' moves go in such different
directions?

\item A point cloud rises in a nearly straight band but with clear scatter. Explain
the difference between the cloud and a linear function fit through it, and state what
the word \emph{model} means in this context.

\item For $g(x)=A\,f(b(x-h))+k$, suppose $A<0$ and $b<0$. Describe both reflections
that occur and name the two axes involved.

\item A cloud opens upward like a bowl with its lowest point near $(2,-3)$. Propose
a parent function and initial values of $A,b,h,k$ to start a fit, and explain each
choice.

\item Open the Function Fit Lab. Generate a cloud, identify the parent family, and
record the four constants you settled on with the resulting fit score. State in one
sentence which constant was hardest to determine and why.
\end{exercises}

\end{document}
